% ============================================================
%  Template Beamer IComp / UFAM (não oficial)
%  Instituto de Computação — IComp · PPGI · UFAM
%  Autor: migvanderlei (https://github.com/migvanderlei)
%
%  Idioma: troque brazil → english no documentclass e no babel
%  Compile: pdflatex (duas vezes) ou Overleaf (pdfLaTeX)
% ============================================================

\documentclass[aspectratio=169, 11pt, brazil]{beamer}

\usepackage[T1]{fontenc}
\usepackage[utf8]{inputenc}
\usepackage[brazil]{babel}
\usepackage{beamerthemeIComp}   % tema (mesma pasta / Overleaf)

% ---- Pacotes extras (opcional) ----
\usepackage{booktabs}
\usepackage{amsmath, amssymb}
\usepackage{tikz}
\usetikzlibrary{arrows.meta, positioning}

% ============================================================
%  METADADOS — preencha aqui
% ============================================================
\title{Título da Apresentação}

\subtitle{Tipo de apresentação — ex.: Exame de Qualificação / Seminário}

\author{Seu Nome Completo}

\institute[IComp · UFAM]{%
  Instituto de Computação — IComp\\
  Universidade Federal do Amazonas — UFAM
}

\orientador{Orientador(a): Prof.\ Dr.\ Nome do Orientador}
\contato{seu.email@icomp.ufam.edu.br}
% Opcional — texto extra no rodapé (após o autor):
% \rodape{IComp}

\date{Manaus, Mês de 2026}

% Títulos muito longos (opcional):
% \tamanhotitulo{large}

% Desativar slides divisores de seção (opcional):
% \divisoresfalse

% ============================================================
\begin{document}
% ============================================================

% ------------------------------------------------------------
%  SLIDE 1 — Capa
% ------------------------------------------------------------
\begin{frame}[noframenumbering, plain]
  \titlepage
\end{frame}

% ------------------------------------------------------------
%  SLIDE 2 — Sumário / Outline
% ------------------------------------------------------------
\begin{frame}{Sumário}
  \tableofcontents
\end{frame}

% ------------------------------------------------------------
\section{Introdução}
% ------------------------------------------------------------

\begin{frame}{Introdução}
  \begin{columns}[T,onlytextwidth]
    \begin{column}{0.52\textwidth}
      \begin{block}{Motivação}
        \begin{itemize}
          \item Contexto e relevância do problema
          \item Desafio principal abordado
          \item Por que este trabalho importa
        \end{itemize}
      \end{block}
    \end{column}
    \begin{column}{0.44\textwidth}
      \begin{destaque}{Contribuição}
        \begin{itemize}
          \item Ideia central em 1--2 frases
          \item Resultado esperado / hipótese
        \end{itemize}
      \end{destaque}
    \end{column}
  \end{columns}
\end{frame}

% ------------------------------------------------------------
\section{Objetivos}
% ------------------------------------------------------------

\begin{frame}{Objetivos}
  \begin{itemize}
    \item Objetivo geral
    \item Objetivo específico 1
    \item Objetivo específico 2
  \end{itemize}
\end{frame}

% ------------------------------------------------------------
\section{Método}
% ------------------------------------------------------------

\begin{frame}{Método}
  \begin{itemize}
    \item Visão geral da abordagem
    \item Etapas principais do pipeline
    \item Premissas e limitações
  \end{itemize}

  \vspace{0.8em}
  \centering
  {\footnotesize\textcolor{icompMidGray}{%
    (Opcional) Insira aqui um diagrama TikZ ou figura.}}
\end{frame}

% ------------------------------------------------------------
%  SLIDE — Figura com legenda (exemplo)
% ------------------------------------------------------------
\begin{frame}{Figura com legenda}
  \centering
  \begin{tikzpicture}[
      box/.style={draw=icompTeal, rounded corners=2pt, fill=icompCanvas,
                  align=center, font=\small, text=icompTealDark,
                  minimum width=2.2cm, minimum height=0.8cm},
      arr/.style={-{Stealth[length=4pt]}, line width=0.8pt, icompTeal}
    ]
    \node[box] (a) {Entrada};
    \node[box, right=1.35cm of a] (b) {Modelo};
    \node[box, right=1.35cm of b] (c) {Saída};
    \draw[arr] (a) -- (b);
    \draw[arr] (b) -- (c);
  \end{tikzpicture}

  \vspace{0.7em}
  {\footnotesize\textcolor{icompDarkGray}{%
    \textbf{Figura 1.} Pipeline ilustrativo: entrada, modelo e saída.}}
\end{frame}

% ------------------------------------------------------------
\section{Resultados}
% ------------------------------------------------------------

\begin{frame}{Resultados}
  \begin{columns}[T,onlytextwidth]
    \begin{column}{0.48\textwidth}
      \centering
      \begin{beamercolorbox}[wd=\textwidth, rounded=true, sep=8pt, center]{block body}
        \vspace{1.4cm}
        {\small\textit{Inserir figura / gráfico aqui}}
        \vspace{1.4cm}
      \end{beamercolorbox}
      \vspace{0.35em}
      {\footnotesize\textit{Fig.~1: Legenda da figura.}}
    \end{column}
    \begin{column}{0.48\textwidth}
      \begin{block}{Observações}
        \begin{itemize}
          \item Achado principal
          \item Comparação com baseline
          \item Interpretação breve
        \end{itemize}
      \end{block}
    \end{column}
  \end{columns}
\end{frame}

\begin{frame}{Tabela de resultados}
  \centering
  \small
  \renewcommand{\arraystretch}{1.2}
  \begin{tabular}{@{}lcc@{}}
    \toprule
    \textbf{Método} & \textbf{Métrica} & \textbf{Valor} \\
    \midrule
    Baseline & --- & --- \\
    Proposto & --- & --- \\
    \bottomrule
  \end{tabular}

  \vspace{0.75em}
  {\footnotesize\textcolor{icompDarkGray}{%
    \textbf{Tabela 1.} Substitua pelos seus números.}}
\end{frame}

% ------------------------------------------------------------
%  SLIDE — Figura ampla (frame plain = mais área útil)
% ------------------------------------------------------------
\begin{frame}[plain]
  \begin{figuraampla}{Figura 2. Substitua pelo seu plot ou \texttt{\textbackslash includegraphics}.}
    \begin{beamercolorbox}[wd=0.92\textwidth, rounded=true, sep=12pt, center]{block body}
      \vspace{1.8cm}
      {\small\textit{Inserir figura aqui}}
      \vspace{1.8cm}
    \end{beamercolorbox}
  \end{figuraampla}
\end{frame}

% ------------------------------------------------------------
\section{Conclusões}
% ------------------------------------------------------------

\begin{frame}{Conclusões}
  \begin{itemize}
    \item \textbf{Conclusão 1:} resultado principal
    \item \textbf{Conclusão 2:} implicação / contribuição
  \end{itemize}

  \vspace{1em}
  \begin{alerta}{Próximos passos}
    Descreva brevemente o trabalho futuro.
  \end{alerta}
\end{frame}

% ------------------------------------------------------------
\section*{Referências}
% ------------------------------------------------------------
\begin{frame}[allowframebreaks]{Referências}
  \begin{thebibliography}{9}
    \bibitem{exemplo}
      Autor~A.
      \newblock Título do trabalho.
      \newblock \emph{Conferência / Periódico}, 2026.
  \end{thebibliography}
\end{frame}

% ------------------------------------------------------------
%  SLIDE FINAL — Agradecimentos
% ------------------------------------------------------------
\textofinal{Obrigado!}
% \rodapefinal{Seu Nome\\seu.email@icomp.ufam.edu.br}
\apoio{%
  \imgapoio[height=0.55cm]{logo_cnpq.png}
  \imgapoio[height=0.55cm]{logo_capes.png}
  \imgapoio[height=0.72cm]{logo_fapeam.png}
}
\agradecimentos

% ============================================================
\end{document}
% ============================================================
